\begin{center}\itshape
A typical integer near $10^{80}$ has about five prime factors, give or take two.\\
The Gaussian law of Erd\H{o}s and Kac is true,\\
but it is a law about numbers nobody will ever write down.
\end{center}

In 1917 G.~H.~Hardy and Srinivasa Ramanujan proved that almost every integer $n$ has about $\log\log n$ distinct prime factors. They also proved an upper bound for the number $\pi_k(x)$ of $n\le x$ with exactly $k$ prime factors, valid for all $k$ at once. It has the right shape, but it does not give the true size of $\pi_k(x)$ when $k$ grows with $x$; Landau had found that size for fixed $k$ in 1900 (an exercise of the Selberg--Delange section). The uniform question is the ``problem of Hardy'' in the title of Sathe's papers of 1953 and 1954, and Selberg's note answered it with the Selberg--Delange method of this chapter.

Between those dates came the theorem of Erd\H{o}s and Kac (1940), told in Volume~5, Chapter~9: $\omega(n)$ is asymptotically Gaussian, with mean and variance $\log\log x$. It is one of the most beautiful theorems in mathematics and one of the least visible. The variance $\log\log x$ grows so slowly that for $x=10^{80}$, about the number of atoms in the observable universe, it is only $5.2$; and at every $x$ the law of $\omega(n)$ is a law on the integers, while the Gaussian is continuous. How far from the Gaussian is the law at a given $x$? In 1949 William LeVeque conjectured that the error is of order $\frac1{\sqrt{\log\log x}}$, the rate of the Berry--Esseen theorem, and in 1958 Alfréd Rényi and Pál Turán proved it, using Selberg's formula. A Hungarian probabilist and a Hungarian number theorist were the right pair for the job.

This chapter draws the consequences of the Selberg--Delange method for $\omega(n)$: the Sathe--Selberg formula for $\pi_k(x)$, uniform in $k\le A\log\log x$; the Rényi--Turán rate in the Erd\H{o}s--Kac theorem; and the large deviations of $\omega(n)$, which are those of a Poisson variable. Throughout,
\[
L=\log x,\qquad\lambda=\log\log x,\qquad x\ge16\ (\text{so that }\lambda\ge1),
\]
and $P_x$, $\mathbb E_x$ refer to $n$ uniform in $\{1,\dots,\lfloor x\rfloor\}$.

\begin{reminder}[from the Selberg--Delange method and Volumes~4, 5 and~10]
\emph{Selberg's formula} (Corollary~\ref{cor:omegasd}). Uniformly for $|z|\le R$, $\sum_{n\le x}z^{\omega(n)}=xL^{z-1}(Q(z)+O_R(\frac1L))$, where $Q(z)=\frac{F(z)}{\Gamma(z)}$ and $F(z)=\prod_p(1+\frac{z-1}p)(1-\frac1p)^{z-1}$. Both $F$ (Lemma~\ref{lem:eulerprod}, as a function of $z$: a locally uniform limit of entire functions) and $Q$ are entire, and $F(0)=F(1)=Q(1)=1$.
\emph{Esseen's smoothing inequality} (Volume~10, Lemma~4.5). If $G$ has a density bounded by $m$, then $\sup|F-G|\le\frac1\pi\int_{-T}^T\frac{|\varphi(t)-\gamma(t)|}{|t|}\dd t+\frac{24m}{\pi T}$.
\emph{Stirling} (Volume~4, Theorem~11.8). $\sqrt{2\pi n}(\frac ne)^n\le n!\le e^{1/12}\sqrt{2\pi n}(\frac ne)^n$ for $n\ge1$.
\emph{Erd\H{o}s--Kac} (Volume~5, Theorem~9.7) and \emph{mod-Poisson convergence} (Volume~10, Theorems~11.6 and~11.7).
\end{reminder}

\section{Large deviations}

The simplest consequence of Selberg's formula uses only real $z$: Markov's inequality applied to $z^{\omega(n)}$ is the Chernoff bound of Volume~3, and the answer is that of a Poisson variable.

\begin{proposition}\label{prop:chernoffomega}
Let $Q^*(\alpha)=\alpha\log\alpha-\alpha+1$ for $\alpha>0$, and fix $A\ge1$. Uniformly for $x\ge16$,
\[
\begin{gathered}P_x\bigl(\omega(n)\ge\alpha\lambda\bigr)\ll_A(\log x)^{-Q^*(\alpha)}\quad(1\le\alpha\le A),\\P_x\bigl(\omega(n)\le\alpha\lambda\bigr)\ll_A(\log x)^{-Q^*(\alpha)}\quad(0<\alpha\le1).\end{gathered}
\]
\end{proposition}
\begin{proof}
For real $0<z\le A$, Selberg's formula with $R=A$ gives $\mathbb E_xz^{\omega(n)}\le C_AL^{z-1}$, since $Q$ is bounded on $[0,A]$ and $\frac x{\lfloor x\rfloor}\le2$. If $\alpha\ge1$, take $z=\alpha$: $z^{\omega(n)}\ge z^{\alpha\lambda}$ on the event, so by Markov's inequality the probability is at most $C_AL^{\alpha-1}\alpha^{-\alpha\lambda}=C_AL^{\alpha-1-\alpha\log\alpha}$, because $\alpha^{\lambda}=e^{\lambda\log\alpha}=L^{\log\alpha}$. If $\alpha\le1$, take $z=\alpha\le1$: then $\omega(n)\le\alpha\lambda$ implies $z^{\omega(n)}\ge z^{\alpha\lambda}$, and the same computation applies.
\end{proof}

For $\alpha=2$ the exponent is $2\log2-1=0.386$: the integers with at least twice the typical number of prime factors are rare, but only by a power of $\log x$. Section~\ref{sec:ld} shows that the exponent is the truth.

\section{The Sathe--Selberg formula}

Let
\[
G(z)=\frac{F(z)}{\Gamma(z+1)}=\frac{Q(z)}z,
\]
an entire function with $G(0)=1$. For $z\ge0$ every factor $1+\frac{z-1}p\ge1-\frac1p$ of $F$ is positive, so $F(z)>0$ and $G(z)>0$ on $[0,\infty[$.

\begin{theorem}[Sathe, 1953; Selberg, 1954]\label{thm:sathe}
Fix $A\ge1$. Uniformly for $x\ge16$ and $1\le k\le A\lambda+1$,
\[
\pi_k(x)=\frac xL\,\frac{\lambda^{k-1}}{(k-1)!}\Bigl(G\Bigl(\frac{k-1}\lambda\Bigr)+O_A\Bigl(\frac k{\lambda^2}\Bigr)\Bigr).
\]
\end{theorem}

The prefactor is worth reading probabilistically: $\frac1L\frac{\lambda^{k-1}}{(k-1)!}=e^{-\lambda}\frac{\lambda^{k-1}}{(k-1)!}$ is the probability that a Poisson variable of mean $\lambda$ equals $k-1$. So $P_x(\omega(n)=k)$ is the probability that $1+\mathrm{Poisson}(\lambda)$ equals $k$, corrected by the factor $G(\frac{k-1}\lambda)$. The ``$1+$'' is the largest prime factor of $n$, which is always there and is not random in number; the Poisson part counts the others.

\begin{proof}
All constants depend only on $A$. Let $S(z)=\sum_{n\le x}z^{\omega(n)}=\sum_k\pi_k(x)z^k$, a polynomial. By Selberg's formula with $R=A+1$, $S(z)=xL^{z-1}(Q(z)+\varepsilon(z))$ with $|\varepsilon(z)|\le\frac{C_1}L$ for $|z|\le A+1$. For every $r>0$, Cauchy's formula gives $\pi_k(x)=M_k+E_k$, where, since $L^z=e^{\lambda z}$ and $\frac{Q(z)}{z^{k+1}}=\frac{G(z)}{z^k}$,
\[
M_k=\frac xL\,\frac1{2\pi i}\oint_{|z|=r}\frac{e^{\lambda z}G(z)}{z^k}\dd z,\qquad|E_k|\le\frac xL\,\frac{C_1}L\,\frac{e^{\lambda r}}{r^k}.
\]
The integral $M_k$ does not depend on $r$.

\emph{The case $k=1$.} $M_1=\frac xLG(0)=\frac xL$. With $r=\frac1\lambda$, $|E_1|\le\frac xL\frac{eC_1\lambda}L$, and $\frac\lambda L=\lambda e^{-\lambda}\le\frac{27e^{-3}}{\lambda^2}$ since $\lambda^3e^{-\lambda}\le27e^{-3}$. This is the claim, with $G(0)=1$.

\emph{The saddle point.} Let $k\ge2$, $y=k-1\ge1$ and $r=\frac y\lambda\in[\frac1\lambda,A]$. The function $e^{\lambda z}z^{-k}$ is smallest on the positive axis at $z=r$ among the circles centred at~$0$ (this is why $r$ is chosen so); expand $G$ around $r$:
\[
G(z)=G(r)+G'(r)(z-r)+(z-r)^2g(z),\qquad g(z)=\frac1{2\pi i}\oint_{|w|=A+1}\frac{G(w)\dd w}{(w-r)^2(w-z)},
\]
the remainder formula of Taylor's theorem for holomorphic functions; for $|z|\le A$, $|g(z)|\le C_2$. Now
\[
\begin{gathered}\frac1{2\pi i}\oint\frac{e^{\lambda z}}{z^k}\dd z=\frac{\lambda^{k-1}}{(k-1)!},\\\frac1{2\pi i}\oint\frac{e^{\lambda z}(z-r)}{z^k}\dd z=\frac{\lambda^{k-2}}{(k-2)!}-r\frac{\lambda^{k-1}}{(k-1)!}=\frac{\lambda^{k-1}}{(k-1)!}\Bigl(\frac{k-1}\lambda-r\Bigr)=0.\end{gathered}
\]
The linear term vanishes: that is the point of the choice of $r$. On $|z|=r$, $z=re^{i\theta}$, we have $|z-r|^2=r^2|e^{i\theta}-1|^2=2r^2(1-\cos\theta)\le r^2\theta^2$ and $|e^{\lambda z}|=e^{y\cos\theta}\le e^{y}e^{-2y\theta^2/\pi^2}$, using $1-\cos\theta\ge\frac{2\theta^2}{\pi^2}$ on $[-\pi,\pi]$. So the remainder term is at most
\[
\frac1{2\pi}\int_{-\pi}^\pi C_2\,r^2\theta^2\,e^ye^{-2y\theta^2/\pi^2}\,r^{1-k}\dd\theta\le C_2r^{3-k}e^y\int_{\R}\theta^2e^{-2y\theta^2/\pi^2}\dd\theta=C_3\,r^{3-k}e^yy^{-3/2}.
\]

\emph{Comparison with the main term.} By Stirling, $(k-1)!=y!\le e^{1/12}\sqrt{2\pi y}\,y^ye^{-y}$, and $r=\frac y\lambda$, so
\[
\frac{(k-1)!}{\lambda^{k-1}}\,C_3r^{3-k}e^yy^{-3/2}\le C_4\,\frac{y^{1/2}y^y}{\lambda^y}\cdot\frac{y^{2-y}}{\lambda^{2-y}}\cdot y^{-3/2}=C_4\frac y{\lambda^2},
\]
and similarly $\frac{(k-1)!}{\lambda^{k-1}}\frac{C_1}L\frac{e^{\lambda r}}{r^k}\le C_5\frac{\lambda}{L}\,y^{-1/2}\le\frac{C_6}{\lambda^2}$. Hence $M_k=\frac xL\frac{\lambda^{k-1}}{(k-1)!}(G(r)+O(\frac k{\lambda^2}))$ and $E_k$ is of the same form.
\end{proof}

\begin{remark}
Volume~10, Theorem~11.7, applied to Selberg's formula (which says that $\omega(n)$ converges mod-Poisson with parameters $\lambda$ and limit $Q$), gives $P_x(\omega(n)=k)=\pi_\lambda(k)(Q(\frac k\lambda)+o(1))$ for $a\le\frac k\lambda\le b$. This is the same statement: $\pi_\lambda(k)Q(\frac k\lambda)=\pi_\lambda(k-1)\frac\lambda kQ(\frac k\lambda)=\pi_\lambda(k-1)G(\frac k\lambda)$, and $G(\frac k\lambda)-G(\frac{k-1}\lambda)=O(\frac1\lambda)$. The saddle-point proof gives the explicit error and covers $k$ near $1$ as well.
\end{remark}

\begin{figure}[ht]
\centering\fig{ch03-local}
\caption{The proportion of $n\le10^8$ with $\omega(n)=k$ (bars), against the Sathe--Selberg main term $e^{-\lambda}\frac{\lambda^{k-1}}{(k-1)!}G(\frac{k-1}\lambda)$ (dots), the law of $1+\mathrm{Poisson}(\lambda)$ (dashed) and the Gaussian of Erd\H{o}s and Kac (grey), with $\lambda=\log\log10^8=2.91$. At $k=3$ the three predictions are $0.274$, $0.230$ and $0.233$, while the truth is $0.348$; the factor $G$ moves Poisson in the right direction, but with $\lambda<3$ the error term $\frac k{\lambda^2}$ is not small.}
\label{fig:sathelocal}
\end{figure}

\section{The rate in the Erd\H{o}s--Kac theorem}

\begin{theorem}[Rényi and Turán, 1958]\label{thm:renyituran}
There is an absolute constant $C$ such that for all $x\ge16$,
\[
\sup_{u\in\R}\Bigl|P_x\Bigl(\frac{\omega(n)-\lambda}{\sqrt\lambda}\le u\Bigr)-\Phi(u)\Bigr|\le\frac C{\sqrt{\log\log x}},
\]
where $\Phi$ is the standard Gaussian distribution function.
\end{theorem}
\begin{proof}
Let $Y=\frac{\omega(n)-\lambda}{\sqrt\lambda}$ and $\varphi(t)=\mathbb E_xe^{itY}$. Define $\varepsilon(z)$ by $\mathbb E_xz^{\omega(n)}=e^{\lambda(z-1)}(Q(z)+\varepsilon(z))$; it is an entire function of $z$, since $\mathbb E_xz^{\omega(n)}$ is a polynomial. By Selberg's formula with $R=2$ and $|\frac x{\lfloor x\rfloor}-1|\le\frac2x\le\frac1L$, $|\varepsilon(z)|\le\frac{C_1}L$ for $|z|\le2$. And $\varepsilon(1)=1-Q(1)=0$. So $\frac{\varepsilon(z)}{z-1}$ is holomorphic, bounded by $\frac{C_1}L$ on the circle $|z-1|=1$, hence on the disc it bounds, by the maximum principle: $|\varepsilon(z)|\le\frac{C_1}L|z-1|$ for $|z-1|\le1$. Similarly $|Q(z)-1|\le C_2|z-1|$ and $|Q(z)|\le C_2$ there.

Let $|t|\le T=\sqrt\lambda$ and $u=\frac t{\sqrt\lambda}$, so $|u|\le1$, and put $z=e^{iu}$, so $|z-1|\le|u|$. Then
\[
\varphi(t)=e^{-it\sqrt\lambda}\,\mathbb E_xz^{\omega(n)}=e^{a}\bigl(Q(z)+\varepsilon(z)\bigr),\qquad a=\lambda(e^{iu}-1-iu).
\]
Let $b=-\frac{t^2}2$. By Taylor's formula $|e^{iu}-1-iu+\frac{u^2}2|\le\frac{|u|^3}6$, so $|a-b|\le\frac{|t|^3}{6\sqrt\lambda}$. Also $\operatorname{Re}a=\lambda(\cos u-1)\le-\frac{2\lambda u^2}{\pi^2}=-\frac{2t^2}{\pi^2}$, and $b\le-\frac{2t^2}{\pi^2}$, so on the segment from $a$ to $b$ the real part stays below $-\frac{2t^2}{\pi^2}$, and $|e^a-e^b|\le|a-b|e^{-2t^2/\pi^2}$. Therefore
\[
|\varphi(t)-e^{-t^2/2}|\le|e^a-e^b|\,|Q+\varepsilon|+e^{-t^2/2}|Q+\varepsilon-1|\le C_3\frac{|t|^3e^{-2t^2/\pi^2}+|t|e^{-t^2/2}}{\sqrt\lambda}.
\]
Esseen's inequality, with the Gaussian density bounded by $m=\frac1{\sqrt{2\pi}}$, gives
\[
\sup_u|P_x(Y\le u)-\Phi(u)|\le\frac{C_3}{\pi\sqrt\lambda}\int_{\R}\bigl(t^2e^{-2t^2/\pi^2}+e^{-t^2/2}\bigr)\dd t+\frac{24}{\pi\sqrt{2\pi}\sqrt\lambda}=O\Bigl(\frac1{\sqrt\lambda}\Bigr).\qedhere
\]
\end{proof}

The rate cannot be improved: $\omega(n)$ is an integer, so the distribution function of $Y$ jumps by $P_x(\omega(n)=k)$ at $u=\frac{k-\lambda}{\sqrt\lambda}$, and for $k$ near $\lambda$ the jump is of order $\frac1{\sqrt\lambda}$ by Theorem~\ref{thm:sathe} (Exercise~\ref{ex:rateoptimal}). A continuous approximation must miss half of it.

\begin{figure}[ht]
\centering\fig{ch03-kolmogorov}
\caption{The Kolmogorov distance $D(x)$ between the law of $\frac{\omega(n)-\lambda}{\sqrt\lambda}$, $n\le x$, and the Gaussian (blue), and $\sqrt{\log\log x}\,D(x)$ (red), for $x$ from $10^3$ to $10^8$. The distance decreases from $0.325$ to $0.239$; the rescaled distance stays near $0.4$, as Theorem~\ref{thm:renyituran} predicts. If it stays there, the distance reaches $0.1$ only around $\lambda\approx17$, that is, $x\approx e^{e^{17}}$.}
\label{fig:kolmogorov}
\end{figure}

\section{Large deviations, sharply}\label{sec:ld}

\begin{corollary}\label{cor:ldomega}
For fixed $\alpha>1$,
\[
\log P_x\bigl(\omega(n)\ge\alpha\lambda\bigr)=-Q^*(\alpha)\log\log x+O_\alpha(\log\log\log x),
\]
and likewise for $P_x(\omega(n)\le\alpha\lambda)$ when $0<\alpha<1$.
\end{corollary}
\begin{proof}
The upper bounds are Proposition~\ref{prop:chernoffomega}. For the lower bound when $\alpha>1$, let $k=\lceil\alpha\lambda\rceil$, so $P_x(\omega(n)\ge\alpha\lambda)\ge\frac{\pi_k(x)}x$. By Theorem~\ref{thm:sathe} with $A=\alpha+1$, and since $G$ is positive and continuous on $[0,A]$, hence bounded below there by some $g_A>0$, we get $\frac{\pi_k(x)}x\ge\frac{g_A}2e^{-\lambda}\frac{\lambda^{k-1}}{(k-1)!}$ for $\lambda$ large. By Stirling, with $k-1=\alpha\lambda+O(1)$,
\[
\log\Bigl(e^{-\lambda}\frac{\lambda^{k-1}}{(k-1)!}\Bigr)=-\lambda+(k-1)\log\lambda-(k-1)\log\frac{k-1}e+O(\log\lambda)=-\lambda Q^*(\alpha)+O(\log\lambda).
\]
For $\alpha<1$ take $k=\lfloor\alpha\lambda\rfloor\ge1$ for $\lambda$ large.
\end{proof}

A Poisson variable of mean $\lambda$ has exactly these large deviations, with rate function $Q^*$ (Volume~10, Chapter~7). The Gaussian of Erd\H{o}s and Kac would predict $e^{-\lambda(\alpha-1)^2/2}$ instead, which is wrong away from $\alpha=1$: the Gaussian describes the centre of the law, the Poisson its tails.

\begin{beauty}
Three layers of the same random integer. Near the mean, $\omega(n)$ looks Gaussian, with an error $\frac1{\sqrt{\log\log x}}$ that no continuous law can beat. At distance $\alpha\log\log x$ it looks Poisson, with Cramér's rate function. And at every scale the arithmetic leaves its fingerprint as a correction factor $G$, an Euler product divided by a Gamma function.
\end{beauty}

\section{Exercises}

\begin{exercise}[$\star$]\label{ex:pi2}
Show that $\pi_2(x)=\frac xL(\lambda+\beta+O(\frac1\lambda))$, where $\beta=G'(0)=\gamma+\sum_p\bigl(\frac1{p-1}+\log(1-\frac1p)\bigr)$ and $\gamma$ is Euler's constant.
\end{exercise}

\begin{exercise}[$\star$]\label{ex:ldnumbers}
Compute the exponents in $P_x(\omega(n)\ge2\lambda)$ and $P_x(\omega(n)\le\frac12\lambda)$. For $x=10^{80}$, how large are these probabilities, up to the factor $(\log x)^{o(1)}$?
\end{exercise}

\begin{exercise}[$\star\star$]\label{ex:rateoptimal}
Show that there is $c>0$ such that, for all large $x$, the supremum in Theorem~\ref{thm:renyituran} is at least $\frac c{\sqrt{\log\log x}}$.
\end{exercise}

\begin{exercise}[$\star\star$]\label{ex:mode}
Show that the most frequent value $k^*(x)$ of $\omega(n)$ for $n\le x$ satisfies $k^*(x)=\lambda+O(1)$. (At $x=10^8$ it is $3$, and $\lambda=2.91$.)
\end{exercise}

\begin{exercise}[$\star\star$]\label{ex:renyituranomega}
Prove Theorem~\ref{thm:renyituran} for $\Omega(n)$, using Corollary~\ref{cor:bigomegasd}.
\end{exercise}

\begin{exercise}[$\star\star\star$]\label{ex:tail}
Let $\alpha>1$ be fixed and $k_0=\lceil\alpha\lambda\rceil$. Show that $P_x(\omega(n)\ge\alpha\lambda)=\frac\alpha{\alpha-1}P_x(\omega(n)=k_0)(1+O_\alpha(\frac1\lambda))$.
\end{exercise}

\section{Hints}

\begin{hint}{ex:pi2}
Theorem~\ref{thm:sathe} with $k=2$, and $G(\frac1\lambda)=1+\frac{G'(0)}\lambda+O(\frac1{\lambda^2})$. For $G'(0)$, differentiate $\log F$ and use $\frac{\Gamma'}\Gamma(1)=-\gamma$.
\end{hint}
\begin{hint}{ex:ldnumbers}
$Q^*(2)=2\log2-1$, $Q^*(\frac12)=1-\frac12-\frac12\log2$; and $\log10^{80}=184.2$.
\end{hint}
\begin{hint}{ex:rateoptimal}
Take $k=\lfloor\lambda\rfloor+1$ and use Theorem~\ref{thm:sathe} and Stirling.
\end{hint}
\begin{hint}{ex:mode}
Compute $\frac{\pi_{k+1}(x)}{\pi_k(x)}$ from Theorem~\ref{thm:sathe} for $1\le k\le3\lambda$, and use Proposition~\ref{prop:chernoffomega} for larger $k$.
\end{hint}
\begin{hint}{ex:renyituranomega}
Only two properties of $Q$ were used: holomorphy on $|z-1|\le1$ and $Q(1)=1$.
\end{hint}
\begin{hint}{ex:tail}
The ratios $\frac{\pi_{k+1}}{\pi_k}$ are $\frac\lambda k(1+O(\frac1\lambda))$; compare the tail with a geometric series, above for all $k\le B\lambda$ and below for $k\le k_0+\lambda^{1/3}$, and use Proposition~\ref{prop:chernoffomega} beyond $B\lambda$.
\end{hint}

\section{Solutions}

\begin{solution}{ex:pi2}
\idea The Sathe--Selberg formula at $k=2$ is evaluated at the point $\frac1\lambda$, close to $0$, so one Taylor term of $G$ gives the second-order term.
\writeup With $k=2$, $\pi_2(x)=\frac xL\lambda(G(\frac1\lambda)+O(\frac1{\lambda^2}))$, and $G(\frac1\lambda)=G(0)+\frac{G'(0)}\lambda+O(\frac1{\lambda^2})$ with $G(0)=1$. So $\pi_2(x)=\frac xL(\lambda+G'(0)+O(\frac1\lambda))$. Now $\frac{G'}G=\frac{F'}F-\frac{\Gamma'}\Gamma(z+1)$, and at $z=0$, $F(0)=1$, $\frac{\Gamma'}\Gamma(1)=-\gamma$, and differentiating $\log F(z)=\sum_p\bigl(\log(1+\frac{z-1}p)+(z-1)\log(1-\frac1p)\bigr)$ term by term (the series of derivatives converges uniformly near $0$, its terms being $O(p^{-2})$) gives $\frac{F'}F(0)=\sum_p\bigl(\frac{1/p}{1-1/p}+\log(1-\frac1p)\bigr)=\sum_p\bigl(\frac1{p-1}+\log(1-\frac1p)\bigr)$. Numerically $\beta=1.0346\ldots$; it is Mertens' constant $0.2615$ plus $\sum_p\frac1{p(p-1)}=0.7731$, the second term counting the $n=p^aq^b$ that are not squarefree.
\end{solution}

\begin{solution}{ex:ldnumbers}
\idea Plug into the rate function, and remember that $\log\log x$ is tiny even for astronomical $x$.
\writeup $Q^*(2)=2\log2-1=0.386$ and $Q^*(\frac12)=\frac12-\frac12\log2=0.153$. So $P_x(\omega(n)\ge2\lambda)=(\log x)^{-0.386+o(1)}$ and $P_x(\omega(n)\le\frac\lambda2)=(\log x)^{-0.153+o(1)}$. For $x=10^{80}$, $\log x=184.2$, and these are about $184^{-0.386}=0.13$ and $184^{-0.153}=0.45$, up to the unknown factor $(\log x)^{o(1)}$, which here is not negligible. Even at the scale of the universe, ``rare'' integers are common.
\end{solution}

\begin{solution}{ex:rateoptimal}
\idea A distribution function on the integers jumps, and a continuous one cannot follow a jump: it misses at least half of it.
\writeup Let $k=\lfloor\lambda\rfloor+1$, so $k-1\le\lambda<k$ and $1\le k\le\lambda+1$. By Theorem~\ref{thm:sathe} with $A=1$, $P_x(\omega(n)=k)\ge\frac{\pi_k(x)}x\ge e^{-\lambda}\frac{\lambda^{k-1}}{(k-1)!}(g_1-\frac{C}{\lambda})$ with $g_1=\min_{[0,1]}G>0$. By Stirling with $m=k-1\in[\lambda-1,\lambda]$, $e^{-\lambda}\frac{\lambda^m}{m!}\ge\frac{c}{\sqrt\lambda}$ (the Poisson probability at its mode). So the jump $J$ of $u\mapsto P_x(Y\le u)$ at $u_0=\frac{k-\lambda}{\sqrt\lambda}$ is at least $\frac{c'}{\sqrt\lambda}$ for $x$ large. If $D$ is the supremum, then $P_x(Y\le u_0)-\Phi(u_0)\le D$ and $\Phi(u_0)-P_x(Y<u_0)\le D$ (take $u\uparrow u_0$ and use the continuity of $\Phi$); adding, $J\le2D$.
\end{solution}

\begin{solution}{ex:mode}
\idea Consecutive frequencies have ratio about $\frac\lambda k$, which crosses $1$ at $k\approx\lambda$.
\writeup For $1\le k\le3\lambda$, Theorem~\ref{thm:sathe} with $A=3$ gives
\[
\frac{\pi_{k+1}(x)}{\pi_k(x)}=\frac\lambda k\cdot\frac{G(\frac k\lambda)+O(\frac k{\lambda^2})}{G(\frac{k-1}\lambda)+O(\frac k{\lambda^2})}=\frac\lambda k\Bigl(1+O\Bigl(\frac1\lambda\Bigr)\Bigr),
\]
since $G$ is bounded below on $[0,3]$ and $G'$ is bounded there, and $\frac k{\lambda^2}\le\frac3\lambda$. There is $C$ such that the ratio is $>1$ for $k\le\lambda-C$ and $<1$ for $\lambda+C\le k\le3\lambda$. For $k\ge3\lambda$, Proposition~\ref{prop:chernoffomega} with $\alpha=3$ bounds $\pi_k(x)\le xC'L^{-Q^*(3)}$ with $Q^*(3)=3\log3-2=1.30>1$, while $\pi_{k}(x)$ at $k=\lfloor\lambda\rfloor$ is of order $\frac x{\sqrt\lambda}$; so the mode is not beyond $3\lambda$. Hence $|k^*-\lambda|\le C+1$.
\end{solution}

\begin{solution}{ex:renyituranomega}
\idea The proof of Theorem~\ref{thm:renyituran} uses Selberg's formula only on the disc $|z-1|\le1$, where Corollary~\ref{cor:bigomegasd} provides the analogue for $\Omega$.
\writeup By Corollary~\ref{cor:bigomegasd}, uniformly for $|z|\le\frac32$, $\sum_{n\le x}z^{\Omega(n)}=xL^{z-1}(Q_\Omega(z)+O(\frac1L))$ with $Q_\Omega=\frac{F_\Omega}\Gamma$, holomorphic on $|z|<2$, and $Q_\Omega(1)=F_\Omega(1)=1$. The disc $|z-1|\le1$ is not inside $|z|\le\frac32$, so run the argument of the proof on the disc $|z-1|\le\frac12$: there $\varepsilon_\Omega(z)$ is entire, $O(\frac1L)$, vanishes at $1$, so $|\varepsilon_\Omega(z)|\le\frac{2C}L|z-1|$; and $|Q_\Omega(z)-1|\le C|z-1|$. Take $T=\frac12\sqrt\lambda$, so that $|u|\le\frac12$ and $|e^{iu}-1|\le\frac12$. Every other step is unchanged, and Esseen's inequality with this $T$ still gives $O(\frac1{\sqrt\lambda})$.
\end{solution}

\begin{solution}{ex:tail}
\idea Beyond $\alpha\lambda$ the frequencies decrease geometrically with ratio about $\frac1\alpha$, so the tail is the first term times $\frac1{1-1/\alpha}$.
\writeup Write $p_k=P_x(\omega(n)=k)$. Fix $B>\alpha$ with $Q^*(B)>Q^*(\alpha)+1$. As in Exercise~\ref{ex:mode}, for $k_0\le k\le B\lambda$, $\frac{p_{k+1}}{p_k}=\frac\lambda k(1+\eta_k)$ with $|\eta_k|\le\frac{C}\lambda$.

\emph{Upper bound.} For $k\ge k_0\ge\alpha\lambda$, $\frac\lambda k\le\frac1\alpha$, so $p_{k_0+j}\le p_{k_0}\rho^j$ with $\rho=\frac{1+C/\lambda}\alpha$ as long as $k_0+j\le B\lambda$. Summing, $\sum_{k_0\le k\le B\lambda}p_k\le p_{k_0}\frac1{1-\rho}=p_{k_0}\frac\alpha{\alpha-1}(1+O(\frac1\lambda))$. Beyond $B\lambda$, Proposition~\ref{prop:chernoffomega} bounds the tail by $C'e^{-\lambda Q^*(B)}$, while $p_{k_0}\ge e^{-\lambda Q^*(\alpha)-O(\log\lambda)}$ by the proof of Corollary~\ref{cor:ldomega}; so this part is $O(\frac{p_{k_0}}\lambda)$.

\emph{Lower bound.} For $0\le i<j\le J=\lfloor\lambda^{1/3}\rfloor$, $\frac\lambda{k_0+i}\ge\frac{\lambda}{\alpha\lambda+1+i}\ge\frac1\alpha(1-\frac{2+2i}{\alpha\lambda})$, so $\frac{p_{k_0+j}}{p_{k_0}}\ge\alpha^{-j}\prod_{i<j}(1-\frac{2+2i}{\alpha\lambda})(1-\frac C\lambda)\ge\alpha^{-j}(1-C''\frac{j^2}\lambda)$. Summing over $j\le J$ gives at least $p_{k_0}(\frac\alpha{\alpha-1}-O(\alpha^{-J})-O(\frac1\lambda))$, since $\sum_jj^2\alpha^{-j}<\infty$; and $\alpha^{-J}=O(\frac1\lambda)$. Combining, $P_x(\omega(n)\ge\alpha\lambda)=\sum_{k\ge k_0}p_k=\frac\alpha{\alpha-1}p_{k_0}(1+O(\frac1\lambda))$.
\end{solution}

\section*{Sources and further reading}

\begin{itemize}
\item \emph{Hardy and Ramanujan.} G.~H.~Hardy and S.~Ramanujan, ``The normal number of prime factors of a number $n$'', \emph{Quart.\ J.\ Math.}\ 48 (1917), 76--92.
\item \emph{Sathe and Selberg.} L.~G.~Sathe, ``On a problem of Hardy on the distribution of integers having a given number of prime factors'', I--IV, \emph{J.\ Indian Math.\ Soc.}\ 17 (1953), 63--141, and 18 (1954), 27--81. A.~Selberg, ``Note on a paper by L.~G.~Sathe'', \emph{J.\ Indian Math.\ Soc.}\ 18 (1954), 83--87.
\item \emph{The rate.} W.~J.~LeVeque, ``On the size of certain number-theoretic functions'', \emph{Trans.\ Amer.\ Math.\ Soc.}\ 66 (1949), 440--463, for the conjecture. A.~Rényi and P.~Turán, ``On a theorem of Erd\H{o}s--Kac'', \emph{Acta Arith.}\ 4 (1958), 71--84.
\item \emph{Mod-Poisson view.} E.~Kowalski and A.~Nikeghbali, ``Mod-Poisson convergence in probability and number theory'', \emph{Int.\ Math.\ Res.\ Not.}\ 2010, 3549--3587. V.~Féray, P.-L.~Méliot and A.~Nikeghbali, \emph{Mod-$\phi$ Convergence}, SpringerBriefs in Probability and Mathematical Statistics, 2016, for precise large deviations in this language.
\item \emph{Textbooks.} G.~Tenenbaum, \emph{Introduction to Analytic and Probabilistic Number Theory}, 3rd ed., AMS, 2015, Chapters~II.6 and~III.4, for the Sathe--Selberg theorem and the Erd\H{o}s--Kac theorem with rate. The atoms in the universe: the usual estimate of $10^{80}$ is an order of magnitude, used here only to compute $\log\log10^{80}=5.2$.
\end{itemize}
